\documentclass[11pt]{article}
\usepackage{amsmath,amsthm,amssymb}
\usepackage[margin=1in]{geometry}
\usepackage{booktabs}

\newtheorem{theorem}{Theorem}[section]
\newtheorem{lemma}[theorem]{Lemma}
\newtheorem{proposition}[theorem]{Proposition}
\newtheorem{corollary}[theorem]{Corollary}
\newtheorem{conjecture}[theorem]{Conjecture}
\theoremstyle{definition}
\newtheorem{definition}[theorem]{Definition}
\theoremstyle{remark}
\newtheorem{remark}[theorem]{Remark}
\newtheorem{example}[theorem]{Example}

\newcommand{\ZZ}{\mathbb{Z}}
\newcommand{\CC}{\mathbb{C}}
\newcommand{\RR}{\mathbb{R}}
\newcommand{\Sym}{\mathfrak{S}}
\newcommand{\im}{\operatorname{im}}
\newcommand{\tr}{\operatorname{tr}}
\newcommand{\rk}{\operatorname{rank}}
\newcommand{\Span}{\operatorname{span}}
\newcommand{\Eplus}[1]{E_{#1}^{(+1)}}
\newcommand{\Eminus}[1]{E_{#1}^{(-1)}}
\newcommand{\kernel}{\operatorname{ker}}

\title{The Left-Two Lemma and the Even-Block Gap\\
in the Rank Isolation Theorem}
\author{Clio}
\date{April 16, 2026}

\begin{document}
\maketitle

\begin{abstract}
We prove the \emph{Left-Two Lemma}: for any representation of~$\Sym_n$
and any $k \ge 4$,
\[
  \Eplus{k-3} \cap \Eplus{k-1} \cap \kernel(P_{k-2}P_k) = \{0\},
\]
where $\Eplus{i}$ denotes the $+1$-eigenspace of~$s_i$ and
$P_i = (1+s_i)/2$. This generalises the double adjacent disjointness
argument that closes the even-block gap at $k=4$ in the Rank Isolation
Theorem, and reduces the remaining gap (even blocks $k \ge 6$) to
showing that the cascade image at the gap step lies in~$\Eplus{k-3}$.
We verify computationally that the cascade image avoids the dangerous
subspace for all irreps of~$\Sym_n$ with $n \le 9$, and identify the
precise obstruction: the intermediate projection~$P_{k-4}$ within
block~$k{-}2$ breaks $\Eplus{k-3}$.
\end{abstract}

%% ====================================================================
\section{Setup and context}\label{sec:setup}
%% ====================================================================

We work in the setting of the Rank Isolation Theorem (see companion
documents \texttt{2026-04-15-rank-isolation.tex} and
\texttt{2026-04-16-frobenius-injectivity.tex}).

\begin{definition}
For a representation $V$ of~$\Sym_n$ and a simple transposition
$s_i = (i,\,i{+}1)$, write
\begin{align*}
  \Eplus{i} &= \{v \in V : s_i v = v\}, \\
  \Eminus{i} &= \{v \in V : s_i v = -v\}, \\
  P_i &= \tfrac{1}{2}(1 + s_i) \colon V \to \Eplus{i}.
\end{align*}
\end{definition}

\begin{lemma}[Adjacent disjointness]
\label{lem:adj-disj}
For any representation of~$\Sym_n$ and adjacent indices $|i - j| = 1$:
\[
  \Eplus{i} \cap \Eminus{j} = \{0\}
  \qquad\text{and}\qquad
  \Eminus{i} \cap \Eplus{j} = \{0\}.
\]
\end{lemma}
\begin{proof}
The subgroup $\langle s_i, s_j \rangle \cong \Sym_3$ satisfies
$(s_i s_j)^3 = 1$. If $s_i v = v$ and $s_j v = -v$, then
$(s_i s_j)^3 v = (s_i s_j)^3 v = (-1)^3 v = -v$, contradicting
$(s_i s_j)^3 = 1$.
\end{proof}

The staircase product $\Pi$ for $w_0 \in \Sym_n$ processes blocks
$B_1, B_2, \ldots, B_{n-1}$ where block~$k$ applies
$P_k, P_{k-1}, \ldots, P_1$ in sequence. The \emph{Rank Isolation}
approach proves $\rk(\Pi|_V) = \dim(V^H)$ (with $H = \langle s_1,
s_3, s_5, \ldots \rangle$) by showing:
\begin{enumerate}
  \item[(a)] \textbf{Within-block cascade} (proved): each step
    $P_{j-1}$ after $P_j$ in a block is injective, by adjacent
    disjointness.
  \item[(b)] \textbf{Odd block starts} (proved): at $P_k$ with $k$
    odd ($s_k \in H$), the contraction argument shows the rank drop
    is exactly the expected one.
  \item[(c)] \textbf{Even block starts}: at $P_k$ with $k$ even
    ($s_k \notin H$), $P_k$ should be injective on the cascade image.
    \textbf{Proved for $k = 2, 4$; gap for $k \ge 6$.}
\end{enumerate}

The even-block injectivity at $k$ reduces (via cascade transport) to
a single step: whether $P_{k-2}$ within block~$k{-}1$'s cascade creates
pure $\Eminus{k}$ vectors from the image in~$\Eplus{k-1}$.

%% ====================================================================
\section{The Left-Two Lemma}\label{sec:left-two}
%% ====================================================================

\begin{definition}[Dangerous subspace]
For even $k \ge 4$, the \emph{dangerous subspace} is
\[
  D_k = \Eplus{k-1} \cap \kernel\!\bigl(P_{k-2} P_k\bigr)
  = \bigl\{w \in \Eplus{k-1} : (1+s_{k-2})(1+s_k)w = 0\bigr\}.
\]
A vector $w \in \Eplus{k-1}$ lies in~$D_k$ if and only if $P_{k-2}$
maps~$w$ into~$\Eminus{k}$.
\end{definition}

\begin{lemma}[Characterisation of $D_k$]
\label{lem:dk-char}
If $w \in D_k$, then the $\Eplus{k}$ component
$w_+ = (w + s_k w)/2$ satisfies $w_+ \in \Eminus{k-2}$.
\end{lemma}
\begin{proof}
Since $s_{k-2}$ and $s_k$ commute ($|k - (k{-}2)| = 2$),
the projection $P_{k-2}$ preserves each $s_k$-eigenspace.
Decompose $w = w_+ + w_-$ with $s_k w_\pm = \pm w_\pm$. Then
\[
  (1 + s_k) w = 2 w_+,
\]
so $(1 + s_{k-2})(1 + s_k) w = 2(1 + s_{k-2}) w_+ = 0$ gives
$w_+ \in \kernel(1 + s_{k-2}) = \Eminus{k-2}$.
\end{proof}

\begin{theorem}[Left-Two Lemma]
\label{thm:left-two}
For any representation of\/ $\Sym_n$ and any $k \ge 4$:
\[
  \Eplus{k-3} \cap \Eplus{k-1} \cap \kernel\!\bigl(P_{k-2}P_k\bigr) = \{0\}.
\]
\end{theorem}

\begin{proof}
Let $w$ lie in the triple intersection. By
Lemma~\ref{lem:dk-char}, $w_+ \in \Eminus{k-2}$.

Since $s_{k-3}$ commutes with $s_k$ ($|k{-}3 - k| = 3 \ge 2$),
the operator $s_{k-3}$ preserves each $s_k$-eigenspace. In particular,
$w \in \Eplus{k-3}$ implies $w_+ \in \Eplus{k-3}$.

Now $w_+ \in \Eplus{k-3} \cap \Eminus{k-2}$. Since
$|(k{-}3) - (k{-}2)| = 1$, Lemma~\ref{lem:adj-disj} gives $w_+ = 0$.

Hence $w = w_- \in \Eplus{k-1} \cap \Eminus{k}$. Since
$|(k{-}1) - k| = 1$, Lemma~\ref{lem:adj-disj} gives $w = 0$.
\end{proof}

\begin{remark}[Relation to $k=4$ proof]
For $k = 4$, the theorem reads $\Eplus{1} \cap \Eplus{3}
\cap \kernel(P_2 P_4) = \{0\}$. At the gap step for block~$4$,
the cascade image is in $\Eplus{3}$ (from~$P_3$) and in $\Eplus{1}$
(from $V^H$ transport, since $s_1 \in H$ and $s_1$ commutes with
$s_3, s_4$). So the Left-Two Lemma directly implies
$P_2$ does not create $\Eminus{4}$ vectors, closing the $k = 4$ gap.
This is exactly the ``double adjacent disjointness'' argument of the
April 15 session.
\end{remark}

%% ====================================================================
\section{Application and obstruction for $k \ge 6$}\label{sec:k6}
%% ====================================================================

\subsection{The general pattern}

For even $k \ge 6$, the Left-Two Lemma would close the gap if the
cascade image at the gap step were in~$\Eplus{k-3}$. The cascade image
at the gap step is $I = P_{k-1}(J)$, where $J$ is the output of
blocks $1, \ldots, k{-}2$.

\begin{proposition}[Eigenspace tracking]
\label{prop:eigenspace}
The cascade image~$I$ at the gap step satisfies:
\begin{enumerate}
  \item[\rm(i)] $I \subset \Eplus{k-1}$ (from the just-applied~$P_{k-1}$).
  \item[\rm(ii)] $I \subset \Eplus{1}$ (since $s_1 \in H$, and $s_1$
    commutes with all generators $s_2, \ldots, s_{k-1}$ used in the
    cascade for $k \ge 5$).
  \item[\rm(iii)] $I \not\subset \Eplus{k-3}$ in general.
\end{enumerate}
\end{proposition}

\begin{proof}
(i) and (ii) follow from the cascade structure: every block ends
with~$P_1$, so $J \subset \Eplus{1}$; and $s_1$ commutes with $s_{k-1}$
(distance $\ge 3$ for $k \ge 5$), so $P_{k-1}$ preserves~$\Eplus{1}$.

For (iii): within block $k{-}2$, the cascade applies $P_{k-2}, P_{k-3},
P_{k-4}, \ldots, P_1$. After $P_{k-3}$: the image is
in~$\Eplus{k-3}$. But $P_{k-4}$ uses $s_{k-4}$, which is adjacent
to $s_{k-3}$ ($|k{-}4 - (k{-}3)| = 1$), so $P_{k-4}$ does not
preserve~$\Eplus{k-3}$. The subsequent steps $P_{k-5}, \ldots, P_1$
and $P_{k-1}$ all commute with $s_{k-3}$ and therefore preserve the
(broken) $s_{k-3}$ eigenspace status. So $\Eplus{k-3}$ is not
recovered.
\end{proof}

\subsection{Computational evidence}

\begin{proposition}[Verified avoidance]
For all irreducible representations $V_\lambda$ of\/~$\Sym_n$ with
$n \le 9$ and all even $k$ with $6 \le k \le n{-}1$:
\[
  I \cap D_k = \{0\}
\]
where $I$ is the cascade image and $D_k$ the dangerous subspace.
In particular, $P_k$ is injective on the cascade image at every even
block start.
\end{proposition}

Combined with parabolic reduction (each even~$k$ reduces to~$\Sym_{k+1}$
irreps) and earlier computational verification through $k = 14$, the
Rank Isolation Theorem holds unconditionally for $n \le 16$.

\subsection{Nature of the obstruction}

The avoidance $I \cap D_k = \{0\}$ holds despite the cascade image
\emph{not} lying in~$\Eplus{k-3}$ and \emph{not} having its
$\Eplus{k}$ components in~$\Eplus{k-2}$. Specifically:
\begin{itemize}
  \item The $\Eminus{k-3}$ component of vectors in~$I$ can be
    substantial (ratios $\|v^{(-)}\|/\|v\|$ up to $0.87$ at $k=8$).
  \item The $\Eminus{k-2}$ component of $w_+ = (v + s_k v)/2$ is
    generically nonzero.
  \item The avoidance is \emph{not} explained by orthogonality ($I$ is
    not in~$D_k^\perp$ within~$\Eplus{k-1}$).
\end{itemize}

The mechanism appears to involve the \emph{global} structure of the
cascade image --- the full history of projections constrains the
image to a subspace that is transverse to~$D_k$ for reasons not
captured by any single eigenspace condition.

%% ====================================================================
\section{Summary of proof architecture}\label{sec:summary}
%% ====================================================================

\begin{center}
\begin{tabular}{llll}
\toprule
\textbf{Component} & \textbf{Status} & \textbf{Method} \\
\midrule
Within-block cascade & Proved & Adjacent disjointness \\
Block 2 start ($P_2$) & Proved & Adjacent disjointness \\
Odd block starts ($P_k$, $k$ odd) & Proved & Contraction argument \\
Even block $k = 4$ ($P_4$) & Proved & Left-Two Lemma with $\Eplus{1}$ \\
Even blocks $k \ge 6$ & \textbf{Gap} & Left-Two Lemma $+$ $\Eplus{k-3}$? \\
\midrule
Unconditional range & $n \le 16$ & Parabolic reduction $+$ computation \\
\bottomrule
\end{tabular}
\end{center}

\noindent The Left-Two Lemma reduces the even-block gap to a single
question:

\begin{conjecture}[Cascade $\Eplus{k-3}$ preservation]
There exists a proof that the cascade image at the even-block gap step
avoids $D_k$ for all $k$, either by:
\begin{itemize}
  \item finding a condition weaker than $\Eplus{k-3}$ that the cascade
    image satisfies and that is sufficient to exclude~$D_k$;
  \item an inductive argument using the proved $k=4$ case as base;
  \item a global argument (e.g., total rank in the regular
    representation) that bypasses block-by-block analysis.
\end{itemize}
\end{conjecture}

%% ====================================================================
\section{Verification}\label{sec:verify}
%% ====================================================================

The Left-Two Lemma was verified computationally for all irreps
of~$\Sym_n$ with $n = 7, 8, 9$. The ``left-two'' intersection
$\Eplus{k-3} \cap \Eplus{k-1} \cap \kernel(P_{k-2}P_k)$ is
$\{0\}$ in every case (75+ cases tested), confirming the algebraic
proof. The ``standard'' intersection
$\Eplus{1} \cap \Eplus{k-1} \cap \kernel(P_{k-2}P_k)$ is
\emph{not} $\{0\}$ for most irreps, confirming that $\Eplus{1}$ is
insufficient.

Scripts: \texttt{even\_block\_gap\_analysis.sage},
\texttt{cascade\_eigenspace\_detail.sage},
\texttt{even\_block\_reverse\_cascade.sage}.

\end{document}
